% OSDI 2026 Paper Template
% The 20th USENIX Symposium on Operating Systems Design and Implementation
% July 13-15, 2026, Seattle, WA, USA
%
% Format: ≤12 pages (excluding references), 8.5"x11", 10pt on 12pt leading,
%         two-column, Times Roman, 7"x9" text block
% Camera-ready: ≤14 pages
%
% Official CFP: https://www.usenix.org/conference/osdi26/call-for-papers
% Template source: https://www.usenix.org/conferences/author-resources/paper-templates
%
% IMPORTANT NOTES:
% - OSDI 2026 has two tracks: Research and Operational Systems
% - For Operational Systems track, title must end with "(Operational Systems)"
% - Max 8 submissions per author
% - Papers should be the right length (not padded to 12 pages)
% - Papers ≤6 pages are unlikely to receive full consideration
% - Use anonymized project/system name (different from arXiv/talks)

\documentclass[letterpaper,twocolumn,10pt]{article}
\usepackage{usenix-2020-09}

% Recommended packages
\usepackage[utf8]{inputenc}
\usepackage{amsmath,amssymb}
\usepackage{graphicx}
\usepackage{booktabs}
\usepackage{hyperref}
\usepackage{url}
\usepackage{xspace}

% Custom commands
\newcommand{\system}{SystemName\xspace}  % Replace with your anonymized system name
\newcommand{\eg}{e.g.,\xspace}
\newcommand{\ie}{i.e.,\xspace}
\newcommand{\etal}{\textit{et al.}\xspace}

\begin{document}

% For submission: use anonymized title and Paper #XXX as author
% For Operational Systems track: add "(Operational Systems)" to title
\title{Your Paper Title Here}
% \title{Your Paper Title Here (Operational Systems)}  % For Operational Systems track

\author{Paper \#XXX}  % Anonymized for submission
% Camera-ready:
% \author{
%   Author One\\
%   \textit{Affiliation One}\\
%   \texttt{email@example.com}
%   \and
%   Author Two\\
%   \textit{Affiliation Two}\\
%   \texttt{email@example.com}
% }

\maketitle

\begin{abstract}
% Write your abstract here.
% Guidelines:
% - State what you achieved
% - Why this is hard and important
% - How you do it
% - What evidence you have
% - Your most remarkable result
Your abstract goes here. Replace this text with a concise summary of your
contribution, approach, and key results.
\end{abstract}

\section{Introduction}
\label{sec:intro}

% Motivate the problem and state your contributions clearly.
% OSDI reviewers look for:
% - Significant problem motivation
% - Interesting and compelling solution
% - Practicality and benefits demonstrated
% - Clear contribution articulation
% - Advances beyond previous work

Your introduction goes here. Clearly state the problem, your approach,
and your contributions.

\section{Background and Motivation}
\label{sec:background}

Provide necessary background and motivation for your work.

\section{Design}
\label{sec:design}

Describe the design of your system.

\section{Implementation}
\label{sec:implementation}

Describe implementation details.

\section{Evaluation}
\label{sec:evaluation}

Present your evaluation results.

\subsection{Experimental Setup}

Describe your experimental setup, workloads, and baselines.

\subsection{Results}

Present and discuss your results.

\section{Related Work}
\label{sec:related}

Discuss related work. Organize methodologically, not paper-by-paper.

\section{Conclusion}
\label{sec:conclusion}

Summarize your contributions and key findings.

% Bibliography
{\footnotesize \bibliographystyle{acm}
\bibliography{references}}

\end{document}
